
% Created 2026-09-05 Sat 13:25
% Intended LaTeX compiler: xelatex
\documentclass[12pt]{article}
\usepackage[utf8]{inputenc}
\usepackage[T1]{fontenc}
\usepackage{graphicx}
\usepackage{longtable}
\usepackage{wrapfig}
\usepackage[margin=1in]{geometry}
\usepackage{rotating}
\usepackage{setspace}
\onehalfspacing
\usepackage{fontspec}
\setmainfont{Libertinus Serif}
\usepackage[normalem]{ulem}
\usepackage{amsmath}
\usepackage{parskip}
\usepackage{amssymb}
\usepackage{capt-of}
\setcounter{secnumdepth}{4} % Tells LaTeX to number down to paragraphs
\setcounter{tocdepth}{4}  
\usepackage{hyperref}
\usepackage[style=oscola, backend=biber]{biblatex}
\addbibresource{/Users/krishnajani/Documents/Libib.bib}
\usepackage{csquotes}
\usepackage[british]{babel}
\setlength{\parskip}{0.5\baselineskip}
\author{Krishna Jani}
\date{\today}
\title{Mobilox Innovations v. Kirusa Software}

\begin{document}
\maketitle

\section*{Facts}
\label{sec:orgc7463e7}
\begin{itemize}
\item The appeal stems from an insolvency petition filed under Section 9 by an operational creditor. The respondent, Kirusa Software, was to provide telephone numbers hat were subsequently polled by Mobilox  Innovations. An NDA was executed between the parties, but the respondent violated it by disclosing on its website that it had worked on the Nach Baliye program run by Star TV.

\item The respondent then issued a demand notice under Section  8 of the IBC, stating that failure to repay the due amount would result in a Sectio 9 petition. By email, the appellant replied that a serious, bona‑fide dispute existed between the parties and that the insolvency petition was merely a pressure tactic. No payment was due until the breach of the NDA was resolved.

\item In 2017 the matter was presented before the adjudicating authority. In light of the pre‑existing dispute, the petition was dismissed. An appeal to the NCLAT was allowed on the following grounds: the adjudicating authority acted mechanically, rejecting the application without examining the issue. The dispute was not only vague but appeared motivated to evade liability for the operational debt. Consequently, the impugned order is set aside and remitted for fresh consideration.
\end{itemize}
\section*{Appellant's Contention}
\label{sec:orgedcfc8d}

\begin{itemize}
\item It  was argued by that the NCLT is obliged to dismiss the application since there was a pre-existing dispute between the parties and the respondent in the present case is a operational creditor.
\item Had he been a financial creditor, the presence or absence of dispute is not material.
\item An interpretation issue was raised before the court, since Section 8(2)(a) states "existence of dispute if any \uline{and} record of pendency of the suit or arbitration proceedings filed" The position was that there is a twofold requirement. One is that there must be a presence of the dispute, and second is that there must be a pending suit or arbitration proceeding. Since in the present case an NDA based dispute was recognized but there was no suit or arbitration proceeding pending, the relevant requisite requirements of Section 8 have not been fulfilled.
\item It was argued that the word "\uline{and} record the pendency of suit" must be read as "or." Since the definition of dispute under 5(6) is inclusive, and its legislative history shows that the original bill had a "means" definition, while the final IBC law has a "includes" definition. Therefore, the moment there is an existence of a dispute, that is a real dispute to be tried and not a sham of frivolous or vexatious dispute, the tribunal is bound to dismiss the application.
\item Once there is a bona fide dispute that is shown by the party. It need not be "ascertained" by the adjudicating authority further.
\end{itemize}
\section*{Respondent's Contentions}
\label{sec:org8162707}
It was argued that Section 5(6) covers three things: existence of the amount of debt, quality of goods or services, or breach of a representation or warranty. Since the issue that was present in this case was that involving an NDA, it did not fall within the ambit of a dispute.
\section*{Court's Reasoning}
\label{sec:orgdb4d3d9}
The court, in deciding this issue, has delved deeply into the history, the legislative background, and the preliminary drafting documents considered in drafting the Code.


\noindent The UNCITRAL Legislative Guide specifically states:

\begin{quote}
Principle among the grounds for denial of the application for technical reasons might be those cases where the debtor is found not to satisfy the commencement standard, where the debt is subject to a legitimate dispute or set off in an amount equal to or greater than the amount of the debt where the proceedings will serve no purpose because, for example, secured debt exceeds the value of assets, and where the debtor has insufficient assets to pay for the insolvency administration, and the law makes no other provision for funding the administration of such estates. In case of a creditor application, it might include those cases where a creditor uses insolvency as an inappropriate substitute for debt enforcement procedures to attempt to enforce a viable business out of the marketplace, or to attempt to obtain preferential payments by coercing the debtor. 
\end{quote}


\noindent Further, an Interim Report on Companies Act, in 2003 stated thus:

\begin{quote}
In order to reinstate the debt enforcement function of the statutory demand test for winding up, if a company fails to pay an \uline{undisputed} debt of prescribed value the creditor should be entitled to a winding up order, irrespective of whether it is insolvent or not. 
\end{quote}


\noindent In the Final Report on Companies Act delivered in 2015 it was held:

\begin{quote}
Operational creditors must present an \uline{undisputed} bill, which may be filed at the registered information utility as the requirement for triggering of IRP.
\end{quote}
\noindent In the ``Notes to Clauses'' of the Insolvency and Bankruptcy Code it is said:
\begin{quote}
Clause eight lays down the procedure for the initiation of the corporate insolvency resolution process by an operational creditor. This procedure differs from the procedure applicable to financial creditors as operational debts tend to be small amounts in and may not be accurately reflected on the records of the IU at all times. \uline{The possibility of disputed debts in relation to operational creditors is also higher in comparison to financial creditors such as banks and financial institutions}. Accordingly, the process for initiation of insolvency resolution process differs from the operational creditor. The corporate debtor has a period of ten days from the receipt of the demand notice to or invoice to inform the operational creditor of the existence of a dispute regarding the debt claim or the repayment of the debt. This ensures that operational creditors whose debt claims are usually smaller are not able to put the corporate debtor into insolvency resolution process prematurely or initiate the process for extraneous considerations. Within fourteen days from the receipt of the application, if the adjudicating authority is satisfied to (a) the existence of a default and (b) the other criteria laid down being met, it shall admit the application. The adjudicating authority is not required to look into any other criteria for admission of the application. It is important that parties are not allowed to abuse the legal process using dealing tactics at the admission stage.
\end{quote}

\noindent The court opines that what is important is the existence of the dispute and/or the arbitration proceeding must be pre-existing. That is, it must exist before the receipt of the demand notice or invoice, as the case may be. The history of the bills which eventually became the IBC indicates that the original definition of dispute has now become an inclusive definition, and the word "bona fide" before suit or arbitration proceedings is deleted. Section 9(5)(ii)(d) must be read with Section 8(2)(a), and a notice of an existing dispute must be received. Therefore, the adjudicating authority, while examining an application under Section 9, will have to determine the following:
\begin{enumerate}
\item Whether there is an operational debt beyond Rs. 1 lakh.
\item Whether the documentary evidence furnished shows that the debt is due and payable and has not been paid.
\item Whether there is an existence of a dispute between the parties, or the record of the pendency of a suit or arbitration proceeding filed before the receipt of the demand notice of the unpaid operational debt in relation to such dispute is filed
\end{enumerate}

\noindent In Innovative Industries v. ICICI Bank, it was held

\begin{quote}
  ``The scheme of Section 7 stands in contrast with the scheme under Section 8, where an operational creditor is, on the occurrence of default, to first deliver a demand notice of the unpaid debt. \uline{Under Section 8, too, the corporate debtor can, within a period of 10 days of the receipt of the demand notice, bring to the notice of the operational creditor the existence of a dispute or the record of the pendency of a suit or arbitration proceedings which is pre-existing (that is, before such notice or invoice was received by the corporate debtor). In the case of a corporate debtor who commits the default of a financial debt, the adjudicating authority has merely to see the records of the information utility or other evidence produced by the financial creditor. It is of no matter that the debt is disputed, so long as the debt is due (that is, payable unless interdicted by some law) or has not yet become due in the sense that it is payable at a future date}.''
\end{quote}

\noindent The corporate debtor has been given 10 days to bring to notice any existing dispute between the parties. The original draft Bill of 2015 stated, "the existence of dispute" alone. The word "and" occurring in Section 8(2) must be read as "or", keeping in mind the legislative intent and the fact that an anomalous situation would arise if it is not read as "or" (if read as "and"). Disputes would only stave off the bankruptcy process if they have already been instituted as a suit or as arbitration proceedings, and not otherwise.

This would lead to great hardship in that a dispute may arise a few days before triggering of the insolvency process, in which case, though a dispute may exist, there is no time to approach either an arbitral tribunal or a court. Similarly, where disputes arise but do not reach the arbitral tribunal within 3 years, such persons would be outside the purview of Section 8(2), leading to bankruptcy proceedings commencing against them. Such an anomaly cannot possibly have been intended by the legislature, nor has it so been intended.  

It is settled that, in certain cases, ``and'' can be read as "or" in order to further the object of the statute. As stated in Stroud's Judicial Dictionary, the word "and" has generally a cumulative sense, but sometimes it is, by force of a context, read as "or."
The court has also looked at two judgments: one of the Australian High Court and the other of the Chancery Division of the UK.

The Australian High Court opines as follows:
\begin{quote}
  In my opinion, the expression ``genuine dispute'' connotes a plausible contention requiring investigation and raises much the same sort of considerations as the "serious question to be tried" criterion which arises on an application for an interlocutory injunction or for an extension or removal of a caveat.
  This does not mean that the court must accept uncritically, as giving rise to a genuine dispute, every statement in an affidavit.
  Something more than mere assertion is required, because if that were not so, then anyone could merely say that it did not owe a debt.
\end{quote}

In Scanhill Pty Limited v Century, the Australian Court held that :

\begin{quote}
  The test to be applied is: \textbf{Whether the court is satisfied that there is a serious question to be tried, that the applicant has an offsetting claim}. Certainly, the court \textbf{will not examine the merits of the dispute} other than to see if there is, in fact, a genuine dispute.
  The notion of a genuine dispute in this context suggests that the court must be satisfied that there is a dispute that is not plainly vexatious or frivolous. \textbf{It must be satisfied that there is a claim that may have some substance}.
\end{quote}

Once the operational creditor has filed an application which is otherwise complete, the adjudicating authority must reject the application under Section 952D if:
- the notice of dispute has been received by the operational creditor
- there is a record of dispute in the IU


All that the authority has to see at this stage is whether there is a plausible contention which requires further investigation and whether the dispute is not a patently feeble legal argument, an assertion, or a fact unsupported by evidence. There is no examination of merits at this stage except to the extent indicated above.


\subsection*{NDA as Dispute}

The second issue that arises for the court's concern is whether a dispute involving the NDA amounts to a dispute under Section 56. On this, the court is clear that the definition of dispute is an inclusive one and uses the word "includes" instead of "means".

Secondly, Section 5(6) only deals with suits or arbitration proceedings which must \textit{relate to one of the three subclauses}. Dispute is said to exist as long as there is a real dispute as to payment between parties. There is a clear bona fide dispute, shown through the communications of the parties, that they constituted a plausible breach of the NDA.

Going by the aforesaid test of existence of a dispute, it is clear that, without going into the merits of the dispute, the appellant has raised a plausible contention requiring further investigation. It may be that the claim eventually does not succeed, but what can be said is that the NCLAT incorrect in characterising the defence as vague or motivated to evade liability.


The decision of the NCLAT was therefore set aside, and the CIRP petition was rejected. 


\end{document}
